Los Angeles County, with 10 million residents, is a coherent political community with unified local government, shared economic integration, and strong geographic identity. Yet for congressional representation, LA County is fragmented across 13 districts drawn by Sacramento, many crossing county lines into neighboring areas. Residents of LA County cannot elect representatives to speak for ``LA County''---they are subsumed into state-designed districts. This fragmentation is not unique to Los Angeles, nor is it the product of Republican gerrymandering. Cook County (Chicago), population 5.3 million, is split across 7 districts by a Democratic Illinois legislature. Harris County (Houston), 4.7 million, is split across 6 districts by a Republican Texas legislature. Both parties prioritize state control over local autonomy.

Congressional redistricting allows state legislatures to fragment large counties for strategic purposes: creating majority-minority districts to satisfy the Voting Rights Act, spreading urban votes to help suburban candidates, protecting incumbents, or pursuing partisan advantage through gerrymandering \citep{cox2016efficient, chen2013unintentional}. The result is that natural political communities---counties with centuries of shared governance, economic interdependence, and civic identity---lose representation coherence. Approximately 21\% of Americans live in counties large enough to span multiple congressional districts, yet these communities are carved into pieces by distant state capitals \citep{putnam2000bowling}.

We propose an alternative: counties with populations exceeding twice the ideal congressional district size should receive autonomous representation via the Huntington-Hill apportionment method. This threshold (1.52 million for 2020) ensures that only counties facing inevitable multi-district fragmentation qualify for autonomy, preserving natural community boundaries while allowing smaller counties to remain in state redistricting pools. Using data from 3,142 U.S. counties across the 2020 census, we demonstrate that county-based representation achieves comparable equality to the current state-based system while eliminating fragmentation for one-fifth of the population.

\textbf{Voting Rights Act implications are substantial and positive.} The 25 qualifying counties could support 197-201 majority-minority districts through county-internal redistricting, compared to 68 under current fragmented boundaries---a nearly threefold increase. This occurs because large urban counties have extremely high minority concentrations (Miami-Dade 87.8\%, Queens 75.2\%, Dallas 71.1\%, Harris County 70.5\%, Brooklyn 66.7\%, Cook County 56.1\%, Los Angeles 72.5\%), and county-based representation eliminates state-level cracking of these populations while providing local control over redistricting. Counties would conduct internal VRA-compliant redistricting using the same methods states currently employ, creating compact majority-minority districts where demographic concentrations warrant. Far from threatening minority representation, county-based representation substantially strengthens it by preventing state legislatures from diluting urban minority voting power through geographic fragmentation.

Critically, this is not a partisan proposal disguised as electoral reform. Our analysis shows that \emph{both} Democratic and Republican state legislatures fragment large counties when they control redistricting: 16 counties in Democratic-controlled states (46 million people) and 7 counties in Republican-controlled states (21 million people) span multiple districts. Democratic states fragment their own urban strongholds (Los Angeles, Chicago, New York City) for various strategic purposes, not just to help Democrats. Republican states fragment their large counties (Harris, Maricopa, Miami-Dade) for similar reasons. The common thread is state control, not party advantage.

We frame county-based representation as fundamentally about \emph{locality}: the principle that natural political communities should not be fragmented for state convenience. This principle transcends current partisan alignments. Though large counties are currently urban and Democratic-leaning---creating a short-term Democratic advantage---the partisan composition of counties is historically contingent on shifting demographic patterns. The principle of local autonomy, once established, outlasts party realignments. We argue by comparison to other precedent-setting reforms (one person, one vote; elimination of literacy tests) that principled justifications can support reforms with asymmetric partisan impacts.

Our analysis proceeds as follows. Section 2 reviews literature on redistricting, gerrymandering, and local governance. Section 3 details our methodology: Huntington-Hill apportionment with a 2.0$\times$ ideal threshold, representation equality metrics, and bipartisan fragmentation analysis. Section 4 presents results across three dimensions: (1) which counties qualify and how fragmentation differs by state control, (2) representation equality compared to the current system, and (3) partisan implications with full disclosure. Section 5 discusses normative arguments for local autonomy, addresses concerns about Voting Rights Act compliance and Democratic advantage, and compares to historical precedents. Section 6 concludes with implications for redistricting reform.
